% Pregunta 437 (2013-noviembre-OEP2013, n.º 13): Robinson Crusoe con shock aditivo f(n)+A.
% Ocio l = T - n (T = 6), consumo c. A escala: f(n) = 1,6*sqrt(n), A = 1,3, u = ln c + ln l.
% E = (4; 2,26), n = 2. E' = (4,63; 3,17), n' = 1,37. La frontera sube en paralelo (misma pendiente
% en cada n): solo efecto renta, sube c y sube el ocio, bajan las horas (b).
% Autor: Víctor Gutiérrez Marcos
\documentclass[tikz,border=6pt]{standalone}
\input{tikz-tcee}
\begin{document}
\begin{tikzpicture}[tcee,x=1.05cm,y=1cm]
  \ejes{$\ell$ (ocio)}{$c$}{8}{6.3}
  % T
  \node[below] at (6,0) {$T$};
  % Fronteras
  \draw[curva1,domain=0:2.4495,samples=100,variable=\t] plot ({6-\t*\t},{1.6*\t});
  \draw[curvanueva,desplazada,domain=0:2.4495,samples=100,variable=\t] plot ({6-\t*\t},{1.6*\t+1.3});
  \node[etiqueta,above right,text=tceemorado] at (0,3.92) {$f(n)$};
  \node[etiqueta,above right,text=tceerojo] at (0,5.22) {$f(n)+A$};
  % Desplazamiento A
  \draw[flecha] (1,3.58) -- node[etiqueta,right,text=tceerojo] {$A$} (1,4.85);
  \draw[flecha] (6,0.05) -- node[etiqueta,right,text=tceerojo] {$A$} (6,1.27);
  % Curvas de indiferencia
  \draw[curva2,domain=2.4:6,samples=80] plot (\x,{9.051/\x}) node[etiqueta,right] {$U_0$};
  \draw[curva2,domain=2.95:6,samples=80] plot (\x,{14.69/\x}) node[etiqueta,right] {$U_1$};
  % Equilibrios
  \proyeccion{(4,2.263)}{}{$c$}
  \proyeccion{(4.633,3.171)}{}{$c'$}
  \node[punto,label={[etiqueta]below left:$E$}] at (4,2.263) {};
  \node[punto,label={[etiqueta]above right:$E'$}] at (4.633,3.171) {};
  % Horas desde T
  \draw[guia,|<->|] (4,-0.75) -- node[below,text=black,font=\small] {$n$} (6,-0.75);
  \draw[guia,|<->|,tceerojo] (4.633,-1.45) -- node[below,font=\small] {$n'<n$} (6,-1.45);
  % Nota
  \node[etiqueta,align=left,font=\footnotesize,anchor=north west] at (3.3,6.2)
    {sube en paralelo: misma pendiente $f'(n)$\\en cada $n$, solo efecto renta};
\end{tikzpicture}
\end{document}
